\section{Conclusões}\label{sec:conclusoes}


Este trabalho apresentou uma abordagem de projeto inteligente para pontes de madeira roliça de estradas vicinais, na qual a modelagem paramétrica dos elementos estruturais foi integralmente acoplada a uma rotina de otimização robusta multiobjetivo em conformidade com a ABNT NBR 7190.
A metodologia foi implementada em uma plataforma computacional de acesso livre e aplicada a uma matriz de vinte configurações, combinando quatro vãos teóricos e as cinco classes de resistência de madeiras de florestas nativas, todas sob o trem tipo TB-240.

A plataforma percorreu as vinte células da matriz sem intervenção manual e sem que nenhuma delas exigisse ajuste individual de parâmetros do algoritmo.
Cada célula consumiu 15~000 avaliações de indivíduos, que correspondem a 75~000 avaliações do modelo mecânico completo, dado que cada indivíduo é submetido a $N_c = 5$ pontos da grade de robustez do diâmetro, totalizando 1{,}5 milhão de avaliações no lote.
Todas as vinte células retornaram fronteiras viáveis com 50 soluções não dominadas, e o hipervolume estabilizou entre as gerações 21 e 153, isto é, com folga de pelo menos 147 gerações em relação às 300 executadas, o que confirma que um único par $N_{gen}$, $N_{pop}$ atende a toda a faixa de vãos e classes estudada.
O custo é modesto: 5{,}6~s por célula em um computador pessoal, ou cerca de 2~min para a matriz completa.
O ponto de maior consequência metodológica é que o dimensionamento passa a ser obtido sem arbitramento inicial de seções: o projetista fornece vão, largura de pista, trem tipo e propriedades do material, e recebe o conjunto completo de soluções de compromisso, em vez de verificar uma seção previamente escolhida por experiência ou analogia.
A matriz de resultados foi ainda condensada em duas ferramentas de anteprojeto complementares (Seções~\ref{sec:abacos} e~\ref{sec:equacao}): os ábacos gráficos de leitura direta e uma equação de pré-dimensionamento $V(L, f_{c0,k}) = a\,L^{b}\,f_{c0,k}^{c}$, ajustada por regressão sobre as vinte células com $R^2 = 0{,}999$ e erro relativo médio de 1{,}3\%, que devolve uma estimativa numérica de volume para qualquer resistência de classe dentro do intervalo estudado, sem exigir leitura gráfica nem nova otimização.


O trabalho reotimizou a célula de referência para uma grade de cinco níveis, $\rho \in \{0\%,\ 5\%,\ 10\%,\ 15\%,\ 20\%\}$ (Seção~\ref{sec:custo_robustez}), para quantificar o custo da robustez, e o resultado desta versão da matriz é positivo e monotônico: a solução de menor volume da fronteira encareceu 6{,}0\% em $\rho = 5\%$, 12{,}8\% em $\rho = 10\%$, 20{,}6\% em $\rho = 15\%$ e 26{,}9\% em $\rho = 20\%$ em relação ao caso determinístico, com taxa marginal levemente crescente entre 1{,}2\% e 1{,}4\% de volume adicional por ponto percentual de $\rho$.
Esse achado contrasta com o de uma versão anterior deste experimento, que não detectava custo de robustez mensurável para $\rho$ até 10\%; a diferença é atribuída à revisão dos limites de busca das variáveis de projeto, que desloca a solução de referência para uma região do espaço de busca mais sensível à perturbação geométrica do diâmetro.
A monotonicidade estrita observada em toda a grade de cinco níveis, ausente na versão anterior, é o que dá a este resultado confiança superior à de uma única execução isolada por nível, ainda que a repetição de cada nível com múltiplas sementes continue sendo o padrão desejável para uma conclusão quantitativa definitiva.
No valor de projeto $\rho = 5\%$ adotado no restante da matriz, o custo de robustez é, portanto, real, mas modesto: cerca de 6\% de acréscimo de volume no extremo econômico, em troca de proteção contra a violação de estados limites sob o desvio dimensional natural da madeira roliça.

A varredura paramétrica produziu três achados.
O primeiro é que o consumo de madeira cresce com o vão segundo uma lei de potência $V \propto L^{n}$ cujo expoente varia pouco entre classes, entre $n = 1{,}64$ (D40) e $n = 1{,}73$ (D20), com queda até a D40 seguida de recuperação parcial nas duas classes seguintes, sem monotonicidade estrita: quanto menor a resistência, mais forte a realimentação entre acréscimo de seção e acréscimo de peso próprio, ainda que o efeito sobre o expoente seja discreto nesta faixa de vãos.
O segundo é que o primeiro degrau de classe é o mais rentável, e com folga: passar de D20 para D30 reduz o consumo em 17 a 21\% em qualquer um dos quatro vãos, de 41\% a 48\% de toda a economia disponível entre D20 e D60, ao passo que os degraus seguintes rendem, individualmente, entre 4 e 12\%.
A redução ao subir de classe é, nesta matriz, aproximadamente uniforme entre os quatro vãos (42\% a 44\% de D20 a D60): nenhuma célula satura no piso dimensional de 20~cm do domínio de busca, e o padrão observado é antes de convergência estrutural uniformemente apertada: em todas as vinte células a solução econômica opera a no máximo 1{,}4\% de folga do limite normativo da verificação governante.

O terceiro achado qualifica a expectativa corrente sem a invalidar por completo, e é o de maior consequência prática.
A flexão governa integralmente a matriz nesta versão, quinze das vinte células pelo tabuleiro e cinco pela longarina, sem nenhuma célula governada pela flecha; a flecha da longarina, contudo, chega extremamente perto do limite em três células de vão e classe mais severos (utilização entre 99{,}2\% e 99{,}97\%), sempre superada pela flexão do tabuleiro por poucos centésimos de ponto percentual.
A propriedade determinante do material permanece, em toda a faixa de 3 a 6~m, a resistência à flexão, com o índice quase permanente de deslocamento abaixo de 39{,}3\% em todas as vinte células; a combinação rara ($L/500$, sem redução por fluência), contudo, já se aproxima do limite nas classes de maior resistência e vãos maiores, um sinal de que a migração de resistência para rigidez começa a se anunciar dentro da própria faixa estudada, e não apenas além dela.
Igualmente contra a intuição, o cisalhamento permanece longe do limite em toda a matriz, com utilização máxima de 60{,}2\%, apesar de ser tradicionalmente citado como crítico em peças roliças curtas, inclusive após a correção, neste trabalho, das envoltórias de momento e cortante para vãos curtos.
E o tabuleiro, que governa quinze das vinte células e a totalidade da fronteira da célula de referência, não é um elemento secundário do dimensionamento: merece a mesma atenção de projeto dedicada às longarinas.

A análise de sensibilidade global reforça essa mesma direção, com uma inversão pontual relevante.
Duas variáveis concentram praticamente toda a variabilidade das funções limite: o diâmetro das peças roliças, responsável por 96\% a 98\% da sensibilidade das três verificações da longarina, e o espaçamento entre longarinas, responsável por 93{,}0\% da sensibilidade da flexão do tabuleiro.
São também as duas variáveis de maior amplitude absoluta ao longo da fronteira (o espaçamento entre longarinas lidera, com 62{,}2~cm, à frente do diâmetro, com 52{,}2~cm), o que as identifica como as verdadeiras variáveis de decisão do problema, uma por subsistema; o coeficiente de variação, por penalizar artificialmente as variáveis de escala pequena, não é o indicador adequado para essa leitura.
A tradução para o controle de execução é direta e seletiva.
O diâmetro exige controle rigoroso na etapa de seleção e classificação das toras, com rejeição das peças fora da faixa nominal, e o espaçamento entre longarinas exige gabarito de posicionamento na montagem.
Já a largura e a altura da prancha do tabuleiro e o espaçamento entre peças do tabuleiro admitem tolerâncias de execução mais folgadas sem prejuízo mensurável à segurança.
Em obra vicinal, na qual o controle de qualidade é necessariamente limitado, saber quais duas dimensões merecem atenção e quais três não a merecem é uma orientação de valor prático imediato.

\subsection*{Limitações e trabalhos futuros}

O escopo deste trabalho restringe-se à superestrutura de pontes de vão único simplesmente apoiadas.
O dimensionamento das ligações, dos encontros, dos pilares e das fundações não foi contemplado e constitui a extensão mais imediata da plataforma.

A quantificação do custo da robustez (Seção~\ref{sec:custo_robustez}) baseou-se em uma única execução do NSGA-II por nível de $\rho$; a monotonicidade estrita observada na grade de cinco níveis atenua, mas não elimina, a preocupação com o ruído de convergência de execuções independentes e isoladas.
Repetir cada nível com múltiplas sementes e reportar médias e desvios permanece o padrão desejável para uma conclusão quantitativa definitiva sobre esse custo, nesta ou em versões posteriores da matriz.

A robustez tratada neste trabalho é de natureza geométrica: perturbam-se as dimensões nominais das peças.
A variabilidade das propriedades mecânicas da madeira e das ações não foi incorporada ao processo de busca.
A conexão com análises de confiabilidade estrutural, por meio de FORM ou de simulação de Monte Carlo, permitiria quantificar o índice de confiabilidade das soluções da fronteira e é o desdobramento de maior potencial científico.

Finalmente, o trabalho restringiu-se ao trem tipo TB-240 e às classes normativas de resistência.
A extensão a trens tipo mais severos e, sobretudo, a substituição das classes normativas por propriedades experimentais de espécies nativas específicas permitiria avaliar em que medida o enquadramento em classe representa adequadamente o material efetivamente disponível na região da obra.
A validação contra o projeto de uma ponte vicinal efetivamente executada permanece como a verificação mais desejável e ainda pendente.
